\documentclass[10pt,a4paper]{article}

% Working manuscript structured as a Scientific Data Data Descriptor.
% The journal does not mandate a visual template at first submission.
% This file is deliberately standalone, as requested for revised LaTeX files.
\usepackage[a4paper,margin=2.35cm]{geometry}
\usepackage[T1]{fontenc}
\usepackage[utf8]{inputenc}
\usepackage{lmodern}
\usepackage{microtype}
\usepackage{graphicx}
\usepackage{booktabs}
\usepackage{tabularx}
\usepackage{array}
\usepackage{xcolor}
\usepackage{enumitem}
\usepackage{amsmath}
\usepackage{hyperref}
\usepackage{cleveref}

\definecolor{natureblue}{HTML}{1F5A7A}
\definecolor{draftamber}{HTML}{8A5A00}
\hypersetup{colorlinks=true,linkcolor=natureblue,citecolor=natureblue,urlcolor=natureblue}
\setlength{\parindent}{0pt}
\setlength{\parskip}{0.55em}
\setlist{nosep,leftmargin=1.5em}
\pagestyle{plain}

\newcommand{\todo}[1]{%
  \par\noindent\fcolorbox{draftamber}{draftamber!7}{%
    \parbox{0.94\linewidth}{\textbf{To be completed after the frozen snapshot.} #1}}\par}
\newcommand{\placeholder}[1]{\textcolor{draftamber}{[\textit{#1}]}}

\title{\textbf{A provenance-aware spatial and spatio-temporal data bank for reproducible statistical research}}
\author{Ghislain Geniaux$^{1}$ \and Johnny D'OLIVEIRA$^{1}$\\[0.5em]
\small $^{1}$\placeholder{affiliation to confirm}\\
\small Corresponding author: \placeholder{name and email to confirm}}
\date{Working manuscript -- 21 September 2026}

\begin{document}
\maketitle

\begin{center}
\fcolorbox{natureblue}{natureblue!6}{\parbox{0.92\linewidth}{%
\textbf{Document status.} Structural LaTeX skeleton for a \emph{Scientific Data}
Data Descriptor. Counts describing the bank, repository identifiers, licences,
author information and data citations must be replaced from the dated frozen
snapshot before submission.}}
\end{center}

\begin{abstract}
Spatial and spatio-temporal datasets are distributed across software packages,
publication supplements and domain repositories, often with incomplete links
between source data, analytical roles, spatial support and published model
specifications. We describe a curated, cross-domain data bank that links source
datasets, local artefacts, documented analytical records and controlled
derivatives. The resource records provenance, response and covariate roles,
geometry, coordinates, time, coordinate reference systems, spatial weights,
published formulas, transformations and readiness for reuse. Its construction
combines structured discovery, full-text evidence extraction, a knowledge graph,
human-readable dataset sheets and deterministic consistency checks. The bank is
intended to support transparent reuse and the design of reproducible empirical
evaluations in spatial statistics. \placeholder{Replace with final repository,
release scope and reuse statement; keep below 170 words.}
\end{abstract}

\section*{Background \& Summary}

Simulation is essential in spatial methodology because it provides known
data-generating mechanisms and permits controlled variation of sample size,
spatial dependence, neighbourhood geometry, noise, heterogeneity and
non-linearity. Yet simulation alone cannot establish whether conclusions extend
across empirical domains, spatial supports, temporal structures and measurement
processes. A reusable empirical resource can complement this evidence when it
preserves the information required to reconstruct each spatial analysis.

General-purpose machine-learning collections demonstrate the value of common
interfaces, machine-readable metadata and a clear separation between datasets,
tasks and computational runs. Domain spatial collections, particularly in
hydrology, show that provenance and variable-level documentation can be handled
rigorously at scale. The remaining need addressed here is a cross-domain layer
that links heterogeneous areal, point, network, raster, panel and
spatio-temporal data to their scientific provenance and published analytical
roles.

The data bank therefore distinguishes a source dataset from its local artefacts,
documented analytical records, controlled derivatives and future benchmark
tasks. It records response and covariate roles, formulas, geometry, time,
coordinate reference systems and spatial weights without treating all catalogued
records as immediately benchmark-ready.

\todo{Insert the final bank contribution and quantitative scope from the frozen
snapshot: number of source datasets, fiches, distributable records, artefacts,
panels, derivatives and mature semi-synthetic records. Cite comparable data
collections and remove any remaining unsupported novelty wording.}

\section*{Methods}

\subsection*{Resource scope and units of description}

The resource uses separate identifiers for source datasets, local artefacts,
documented dataset sheets and derived analytical records. A source dataset is
counted once even when several targets, formulas, periods or subsamples are
derived from it. Panel parents remain distinct from cross-sectional extracts.
Prediction products, coefficient surfaces and other model outputs are not
classified as raw observational datasets unless the distributed record is
explicitly designed around that object.

\todo{Freeze the operational definitions of source dataset, artefact, dataset
sheet, panel parent, controlled subsample, semi-synthetic record and benchmark
task.}

\subsection*{Input data discovery}

Datasets were identified through three source families: datasets distributed by
R or Python packages; open datasets linked to scientific publications in
spatial statistics, spatial econometrics and spatio-temporal modelling; and
datasets obtained from public data repositories and portals. For publication-
linked records, full texts were converted to structured TEI, and dataset names,
repository identifiers, empirical roles and model evidence were checked against
the source article. Repository relations and licences were verified before
download. A bibliographic DOI was never treated as a dataset DOI without an
explicit repository relation.

\todo{Report exact searches, dates, versions, inclusion rules and persistent
identifiers for every input collection. Cite each deposited input dataset using
the journal's data-citation format.}

\subsection*{Provenance and curation workflow}

The working workflow proceeds from immutable raw sources to a curated corpus,
structured knowledge graph and human-readable dataset sheets. Automated
extraction produces candidate evidence rather than final assertions. Curators
verify source identity, variables, spatial support, formula evidence,
transformations and limitations. Deterministic audits then compare the knowledge
graph, dataset sheets, file registry and local artefacts.

\subsection*{Spatial and temporal harmonisation}

Where supported by the source, tabular observations are converted to simple-
feature objects while retaining source identifiers and provenance. Coordinates,
geometry and coordinate reference systems are recorded separately from model
covariates. Temporal variables and stable unit identifiers are retained for
panels. Spatial weights supplied by a source are preserved; geographic weights
may be reconstructed only when the geometry and construction rule are documented.

\todo{Insert the final conversion rules, software versions, handling of missing
values, CRS policy, geometry repair policy and checks for panel identifiers.}

\subsection*{Controlled subsampling}

Controlled subsamples remain linked to their parent dataset and record the
selection rule, random seed, retained units and intended purpose. They are not
counted as independent source datasets.

\todo{Describe only the subsamples included in the frozen release.}

\subsection*{Semi-synthetic records}

Semi-synthetic records combine observed spatial support or covariates with an
explicitly documented outcome-generating mechanism. Each record must distinguish
observed inputs, fitted components, controlled interventions and generated
quantities, and must retain the random seed and generating parameters.

\todo{Include only sufficiently mature records. Detailed estimator comparisons
and package benchmarking remain outside this Data Descriptor.}

\section*{Data Records}

The deposited release will be described at file level, including repository,
accession, folder structure, formats, data dictionaries and relationships between
parents and derivatives. Table~\ref{tab:records} is the required summary table;
its values must be generated from the frozen snapshot.

\begin{center}
\refstepcounter{table}\label{tab:records}
\textbf{Table~\thetable.} Planned record groups in the deposited release.
Replace placeholders with frozen counts and repository locations.\par\smallskip
\begin{tabularx}{\textwidth}{>{\bfseries}p{0.22\textwidth}X p{0.18\textwidth}}
\toprule
Record group & Contents & Frozen count \\
\midrule
Metadata registry & Machine-readable identifiers, provenance, status and links & \placeholder{pending} \\
Dataset sheets & Human-readable documentation and evidence & \placeholder{pending} \\
Final spatial artefacts & Reusable spatial or spatio-temporal objects & \placeholder{pending} \\
Panel parents & Longitudinal records with unit and time identifiers & \placeholder{pending} \\
Controlled derivatives & Reproducible subsamples linked to parents & \placeholder{pending} \\
Semi-synthetic records & Generated outcomes and complete generating metadata & \placeholder{pending} \\
\bottomrule
\end{tabularx}
\end{center}

\todo{List the actual repository accession, exact filenames, formats, folder
hierarchy, checksums, schemas and data dictionaries. Do not report summary
statistics here.}

\section*{Data Overview}

\todo{Optional. Limit to one paragraph and at most one or two descriptive
figures or tables. Remove this section if the overview can be generated directly
from the deposited files without aiding reuse.}

\section*{Technical Validation}

Technical validation will report checks on schema conformance, identifier
alignment, file readability, dimensions, variable presence, geometry validity,
coordinate reference systems, temporal indexing, formula consistency,
neighbourhood reconstruction, parent--derivative links and provenance. A record
is not reported as reusable merely because a file was downloaded or a dataset
sheet exists.

\begin{center}
\refstepcounter{table}\label{tab:validation}
\textbf{Table~\thetable.} Validation dimensions to be populated from the frozen audit.\par\smallskip
\begin{tabularx}{\textwidth}{>{\bfseries}p{0.27\textwidth}X p{0.18\textwidth}}
\toprule
Validation dimension & Acceptance evidence & Frozen result \\
\midrule
File integrity & Readable artefact, checksum and documented format & \placeholder{pending} \\
Variable consistency & Declared response and covariates occur in the artefact & \placeholder{pending} \\
Spatial support & Valid geometry or coordinates and documented CRS & \placeholder{pending} \\
Spatial weights & Original or reproducible construction rule & \placeholder{pending} \\
Temporal structure & Valid time field and stable panel identifier where required & \placeholder{pending} \\
Provenance & Source, version, licence and transformation chain recorded & \placeholder{pending} \\
Cross-layer agreement & KG, sheets, registry and files agree & \placeholder{pending} \\
\bottomrule
\end{tabularx}
\end{center}

\todo{Insert reproducible audit commands, sample sizes, failure counts and the
handling of unresolved records. Validation concerns data quality and consistency,
not estimator rankings.}

\section*{Usage Notes}

Users should distinguish catalogued records from distributable artefacts and
from records explicitly ready for a given analytical task. Published formulas,
targets and spatial weights should be reused only with their recorded evidence
and preprocessing. Panel parents should not be silently collapsed into a cross-
section, and random partitions should not be assumed appropriate for spatial or
temporal validation.

\todo{Add loading examples, licence-specific restrictions, software versions,
known limitations and instructions for citing both this bank and original data
creators.}

\section*{Data Availability}

\todo{Repeat the final repository name, DOI or accession, version, principal
record groups and access restrictions. The deposited data must be available to
reviewers at initial submission and public in an appropriate repository before
publication.}

\section*{Code Availability}

\todo{Provide the archived code release, DOI, version or commit, licence and the
commands required to regenerate the registry, audit tables and figures.}

\section*{Acknowledgements}

\placeholder{Funding, institutional and individual acknowledgements.}

\section*{Author Contributions}

\placeholder{Use the CRediT taxonomy and confirm contributions with all authors.}

\section*{Competing Interests}

The authors declare \placeholder{statement to confirm}.

\section*{Funding}

\placeholder{Funding bodies and grant numbers.}

% Embedded bibliography keeps this .tex standalone for revised submission.
% Replace these scaffolding entries with the verified master bibliography.
\begin{thebibliography}{99}
\bibitem{scientificdata-guidelines}
Scientific Data. Submission guidelines. Nature Portfolio
(\url{https://www.nature.com/sdata/publish/submission-guidelines}).

\bibitem{scientificdata-policies}
Scientific Data. Data policies. Nature Portfolio
(\url{https://www.nature.com/sdata/policies/data-policies}).
\end{thebibliography}

\end{document}
